En esta sección se describen las herramientas tecnológicas empleadas para el desarrollo del trabajo terminal, abarcando desde la administración y gestión del proyecto hasta los entornos de codificación y colaboración del equipo.

\subsection{Herramientas de Administración}

Estas herramientas se emplean en las actividades de documentación y gestión del trabajo, conforme a la metodología descrita en el capítulo \ref{cap:marco_metodologico}.

\subsubsection{Editor de Texto}
Toda la redacción técnica relacionada con el proyecto se documenta por medio de \textbf{Overleaf}, un editor de LaTeX disponible directamente en el navegador web.

La principal ventaja de esta herramienta es la posibilidad de trabajar de forma colaborativa en línea, sin necesidad de instalar programas específicos. Asimismo, conserva un historial de versiones del documento y compila en línea, lo que facilita la revisión de cambios entre los integrantes. A diferencia de procesadores de texto como Microsoft Word, maneja documentos extensos de forma consistente, ya que la estructura y el formato se controlan mediante código.

\subsubsection{Software de Gestión del Trabajo}
El control de actividades se lleva a cabo mediante la aplicación de productividad \textbf{Notion}, que permite adaptar un espacio de trabajo compartido y personalizado para todos los integrantes del equipo, con elementos visuales que facilitan el seguimiento de las tareas. A través de esta plataforma se implementa el tablero Kanban para la asignación, gestión y control de actividades. Asimismo, se registran los \textit{sprints} completados, junto con la evidencia de las reuniones del equipo (retroalimentaciones, puntos importantes y grabaciones de las sesiones) o algún otro material de apoyo.

\subsubsection{Plataforma de Videoconferencias}
Todas las reuniones se realizan en \textbf{Google Meet}, una herramienta de fácil acceso, sólida y que no requiere descargas obligatorias. Además, permite grabar las sesiones y generar notas de las reuniones.

\subsection{Herramientas de Desarrollo}
Estas herramientas se emplean en el desarrollo del prototipo de aplicación móvil.

\subsubsection{IDE}
El desarrollo del sistema se realiza en \textbf{Visual Studio Code}, empleado para la escritura, depuración y ejecución del código del prototipo. Se eligió por su versatilidad y por su facilidad de instalación y configuración.

\subsubsection{Framework}
Se utiliza \textbf{Flutter}, un framework de código abierto para el desarrollo de aplicaciones multiplataforma con el lenguaje de programación Dart\cite{flutter_businessvalue}, para construir el cliente móvil del prototipo. Su elección se justifica en la sección \ref{sec:factibilidad_técnica}.

\subsubsection{Sistema Operativo}
El desarrollo se lleva a cabo sobre \textbf{Windows 11}, porque el equipo ya cuenta con este sistema y funciona correctamente con el resto de las herramientas descritas en esta sección.

\subsubsection{Sistema de Control de Versiones}
Se utiliza \textbf{Git} para llevar el historial de cambios del código y facilitar el trabajo colaborativo.

\subsubsection{Gestor de Repositorios}
Se utiliza \textbf{GitHub} para el alojamiento del código y la revisión de los cambios del equipo.

\subsubsection{Gestor de Base de Datos}
Se utiliza \textbf{Supabase}, que provee una base de datos PostgreSQL con soporte para las extensiones PostGIS, para el análisis geoespacial, y pgRouting, para el cálculo de rutas. Además, expone una capa de API sobre la base de datos y cuenta con un paquete para Flutter disponible en pub.dev, lo que facilita la integración con el cliente móvil. La justificación completa se presenta en la sección \ref{sec:factibilidad_tecnica}.

\subsubsection{APIs}
Se utilizan la API de \textbf{Google Street View}, para la obtención de imágenes del entorno, y \textbf{Google Places API}, para la consulta de puntos de interés. Google ofrece una cobertura amplia y actualizada de la zona de estudio, por lo que se eligieron estas APIs como fuentes de datos.

El Cuadro \ref{tab:herramientas_proyecto} resume el software seleccionado para cubrir las características requeridas por el sistema.

\begin{table}[hbt!]
\centering
\renewcommand{\arraystretch}{1.5}
\begin{tabular}{| >{\centering\arraybackslash}m{6cm} | >{\centering\arraybackslash}p{6cm} |}
\hline
\textbf{Características} & \textbf{Herramientas de software} \\ \hline
\textbf{Sistema Operativo} & Windows 11 \\ \hline
\textbf{Editor de Texto} & Overleaf \\ \hline
\textbf{Gestión del trabajo} & Notion \\ \hline
\textbf{Videoconferencias} & Google Meet \\ \hline
\textbf{IDE} & Visual Studio Code \\ \hline
\textbf{Framework} & Flutter (Dart) \\ \hline
\textbf{Sistema de control de versiones} & Git \\ \hline
\textbf{Gestor de repositorios} & GitHub \\ \hline
\textbf{Gestor de Base de Datos} & Supabase (PostgreSQL, PostGIS, pgRouting) \\ \hline
\textbf{APIs} & Google Street View, Google Places API \\ \hline
\end{tabular}
\caption{Herramientas para el desarrollo del proyecto}
\label{tab:herramientas_proyecto}
\end{table}